% SHARD-ID: AQM-05-TORIC-SUPERCHARGE
% SHARD-TITLE: Toric Code Supercharge From Local Checks
% SHARD-SUMMARY: Answers the first toric-code question with an auxiliary-fermion construction.
% SHARD-SUMMARY: Defines the degree-zero Q-cohomology and boundary quotient in elementary terms.
% SHARD-SUMMARY: Unifies the ghost-supercharge construction with the face-edge-vertex boundary complex.
% SHARD-KEYWORDS: toric code, supercharge, ghosts, fermions, local checks, cohomology, homology

\section{Toric Code Supercharge From Local Checks}

\subsection{Question}

Can one formulate the Kitaev toric code concretely as
\[
  H = Q Q^* + Q^* Q
\]
for a supercharge \(Q\), and can the code space be realized by the elementary
``\(Q\)-cohomology'' construction?

The answer recorded here is:

\begin{quote}
Yes, after adding one auxiliary fermion, or ghost, for each local stabilizer
check. This gives a local-check supersymmetric package. It is not yet a
purely physical-qubit supercharge.
\end{quote}

The toric-code stabilizer convention is sourced from the local TeX file
\path{references/toric_code/kitaev_quant_ph_9707021_source/anyons.tex}. Lines
189--222 define qubits on edges, star operators \(A_s\), plaquette operators
\(B_p\), and the protected subspace; lines 328--337 define Kitaev's
Hamiltonian and state that its ground state is the protected subspace.

\subsection{Elementary Construction}

Let \(V\) be the physical Hilbert space of the toric-code qubits. For every
star or plaquette stabilizer \(S_i\), define the violated-check projector
\[
  P_i = \frac{1-S_i}{2}.
\]
Thus \(P_i v=0\) means that the state \(v\) satisfies check \(i\), while
\(P_i v=v\) means that this component violates check \(i\). The shifted
positive toric-code Hamiltonian is
\[
  H_{\mathrm{TC}} = \sum_i P_i .
\]

Now add one auxiliary fermion for each check. Concretely, this means we allow
formal labels \(e_i\), one per check, and impose the usual rule that swapping
two labels changes the sign. Let \(c_i^*\) be the operation ``attach the label
\(e_i\)'', and let \(c_i\) be the operation ``remove the label \(e_i\)'',
with zero result if the label is absent. These operations satisfy
\[
  c_i c_j^* + c_j^* c_i = \delta_{ij}, \qquad
  c_i c_j + c_j c_i = 0, \qquad
  c_i^* c_j^* + c_j^* c_i^* = 0 .
\]

Define the cochain supercharge
\[
  Q = \sum_i c_i^* P_i .
\]
This operator sends a physical state to the formal list of its violated local
checks. Because the projectors \(P_i\) commute, and because two fermion
creations anticommute, \(Q^2=0\). The adjoint \(Q^*=\sum_i c_i P_i\) removes
one ghost label. The cross-terms in \(Q Q^*+Q^* Q\) cancel in pairs, again
using commutation of the projectors and anticommutation of the fermions, so
\[
  Q Q^* + Q^* Q = \sum_i P_i = H_{\mathrm{TC}}
\]
on the physical ghost-vacuum sector, and \(H_{\mathrm{TC}}\otimes I\) on the
full ghost-extended space.

\subsection{What The Cohomology Means}

Here is the promised elementary definition. Degree zero means ``no ghost label
attached'', so a degree-zero state is just a physical toric-code state
\(v\in V\). The degree-zero \(Q\)-cohomology is:

\begin{itemize}
  \item closed degree-zero states: those \(v\) with \(Qv=0\);
  \item exact degree-zero states: those of the form \(Qw\) where \(w\) has
  degree \(-1\).
\end{itemize}

There are no degree \(-1\) states in this construction, so no nonzero exact
degree-zero states. Also, \(Qv=\sum_i e_i\otimes P_i v\), and the one-ghost
labels \(e_i\) are independent. Therefore
\[
  Qv=0
  \quad\Longleftrightarrow\quad
  P_i v=0\ \text{for every check }i
  \quad\Longleftrightarrow\quad
  S_i v=v\ \text{for every check }i .
\]
Thus the degree-zero \(Q\)-cohomology is exactly the toric-code code space.
The qualifier ``degree-zero'' is essential: on a code state all \(P_i\) vanish,
so the full ghost-extended cohomology also contains ghost-labelled copies of
the code space in higher degrees. The physical code is the no-ghost part.

Equivalently, using the boundary operator \(Q^*=\sum_i c_iP_i\), degree-zero
homology is the quotient of all physical states by the span of states with at
least one local-check violation. Since the \(P_i\) are commuting orthogonal
projectors, this quotient has the protected subspace as its canonical
representative.

\subsection{Algebraic Validation}

The validation script is
\path{scripts/lattice_codes/toric_supercharge_validation.jl}. It checks the
stabilizer data algebraically, not by building the full physical or ghost
Hilbert matrices. For the \(k=4\) torus it verifies:

\begin{itemize}
  \item \(n=2k^2=32\) physical qubits and \(32\) local checks;
  \item each check has weight \(4\);
  \item the binary symplectic commutator of every pair of checks is zero;
  \item the stabilizer rank is \(2k^2-2=30\);
  \item the code dimension is \(2^{32-30}=4\);
  \item the stored \texttt{q\_square\_certificate} and
  \texttt{anticommutator\_certificate} booleans are true because the checked
  commuting-projector and involutive-check hypotheses needed by the
  CAR/projector proof hold;
  \item the degree-zero \(Q\)-cohomology dimension is \(4\).
\end{itemize}

The generated CSV is
\path{runs/2026-05-24-toric-supercharge/data/toric_supercharge_summary.csv}.
Its first row is a sentinel noting
that no full physical or ghost Hilbert matrices are built. Thus the run records
the checked hypotheses and summary booleans used by the formal proof above; it
is not a stored matrix-level certificate for \(Q^2=0\) or the anticommutator.

\subsection{Boundary Maps and the Same Construction}

There is a second, more topological description of the toric code. It starts
from the cellular chain complex over \(\mathbb{F}_2\)
\[
  C_2 \xrightarrow{\partial_2} C_1 \xrightarrow{\partial_1} C_0 ,
\]
where \(C_2\) is spanned by faces, \(C_1\) by edges, and \(C_0\) by vertices.
The local Bombin--Martin-Delgado TeX source listed in
\path{references/toric_code/SOURCES.md} defines these maps for 2-complexes in
lines 2050--2140 and records \(\partial^2=0\). The local QECC Zoo source in the
same manifest records the toric-code family as a CSS code with star and
plaquette stabilizers in lines 20--39.

The cellular description does suggest a different supercharge first. Define
the cochain maps
\[
  \delta^0=\partial_1^T:C^0\to C^1,\qquad
  \delta^1=\partial_2^T:C^1\to C^2 .
\]
On the finite graded vector space \(C^0\oplus C^1\oplus C^2\), the cellular
supercharge
\[
  Q_{\mathrm{cell}}=\delta^0+\delta^1
\]
is nilpotent because
\[
  Q_{\mathrm{cell}}^2=\delta^1\delta^0
  =(\partial_1\partial_2)^T=0 .
\]
Its middle cohomology means the following elementary quotient: take all edge
bit strings \(a\in C^1\) with \(\delta^1 a=0\), and identify two such strings
when they differ by \(\delta^0 b\) for some vertex bit string \(b\in C^0\).
In symbols,
\[
  H^1_{\mathrm{cell}}=\ker(\delta^1)/\operatorname{im}(\delta^0).
\]

The quantum CSS unification is then:
\[
  H_X = \partial_1,\qquad H_Z=\partial_2^T .
\]
The rows of \(H_X\) are the vertex-star supports for the \(X\)-checks. The rows
of \(H_Z\), equivalently the columns of \(\partial_2\), are the face-boundary
supports for the \(Z\)-checks. The CSS commutation condition is not extra:
\[
  H_X H_Z^T = \partial_1\partial_2 = 0 .
\]

The precise chain-to-ghost statement is proved in
Section~\ref{sec:chain-ghost-proof}. It gives an explicit map from cellular
cohomology classes to toric-code states and derives the ghost
anticommutator from the CAR relations and the commuting stabilizer projectors.
The operator-level comparison of \(Q_{\mathrm{cell}}\) and \(Q_{\mathrm{gh}}\)
is in Section~\ref{sec:operator-relation}.

Concretely, the boundary map
\(\partial_2:C_2\to C_1\) and \(\partial_1:C_1\to C_0\) act on classical
\(\mathbb{F}_2\)-chains: formal sums of faces, edges, and vertices. They
classify error strings, check supports, and logical operators. The
boundary data determine the stabilizer checks \(S_i\), and those checks
determine the projectors \(P_i=(1-S_i)/2\) used in \(Q\).

The chain complex also explains the number of encoded qubits. The first
homology group is
\[
  H_1=\ker(\partial_1)/\operatorname{im}(\partial_2).
\]
For the torus this has dimension \(2\) over \(\mathbb{F}_2\), so the code
space has dimension \(2^2=4\). In the \(Z\)-basis one sees the dual statement:
the plaquette constraints select \(1\)-cocycles, and the star operators
identify cocycles differing by a coboundary. Hence the code basis is naturally
labelled by \(H^1\), while the logical string operators are labelled by the
paired \(H_1/H^1\) data.

The validation script is
\path{scripts/lattice_codes/toric_chain_ghost_unification.jl}. For \(k=4\), it
verifies \(\partial_1\partial_2=0\), checks that rows of
\(\partial_1\) match the star supports and columns of \(\partial_2\) match the
plaquette supports, and computes
\[
  \operatorname{rank}\partial_1=15,\quad
  \operatorname{rank}\partial_2=15,\quad
  \dim H_1 = 32-15-15=2 .
\]
It also records that \(Q_{\mathrm{cell}}^2=0\), that the middle cellular
cohomology has \(2\) binary generators, that the \(Z\)-basis cocycle space has
\(2^{17}=131072\) elements, and that quotienting by the \(2^{15}=32768\)
star-generated coboundaries leaves \(4\) classes.
The run output is
\path{runs/2026-05-24-toric-chain-ghost-unification/data/toric_chain_ghost_unification.csv}.
